\section{Introduction}

Deep Neural Networks (DNN) have powered a wide spectrum of applications, ranging from image recognition~\cite{resnet}, language translation~\cite{bert}, anomaly detection~\cite{du2017deeplog}, content recommendation~\cite{van2013deep}, to drug discovery~\cite{ramsundar2019deep}, art generation~\cite{mao2017deepart}, game play~\cite{guo2014deep}, and self-driving cars~\cite{bojarski2016end}. Many applications pursue higher intelligence by optimizing larger models using larger datasets, craving advances in distributed training systems. Among existing solutions, distributed data parallel is a dominant strategy due to its minimally intrusive nature. This paper presents the design, implementation, and evaluation of the distributed data parallel package in PyTorch v1.5~\cite{pytorch:nips:2019}.

Training a DNN model usually repeatedly conducts three steps~\cite{lecun1988theoretical}, the forward pass to compute loss, the backward pass to compute gradients, and the optimizer step to update parameters. The concept of data parallelism is universally applicable to such frameworks. Applications can create multiple replicas of a model, with each model replica working on a portion of training data and performing the forward and backward passes independently. After that, model replicas can synchronize either their gradients or updated parameters depending on the algorithm. It's nominally possible to build a working version of data parallel purely on the application side, as it only requires inserting appropriate communications into every iteration. However, squeezing out the last bit of performance takes an enormous amount of effort in design and tuning. Providing native distributed data parallel APIs on the platform side would help application developers focus on optimizing their models, while the platform developing team could continuously and transparently improve the training speed. To provide a general distributed data parallel package, the challenges are three-fold.
\begin{itemize}
    \item \textbf{Mathematical equivalence}: The purpose of data parallel is to speed up training on large datasets. Applications expect to harvest the same result model as if all training had been performed locally without model replication. This requires mathematical equivalence to local training despite its distributed nature.
    \item \textbf{Non-intrusive and interceptive API}: Application developments usually start from local models and then scale out when necessary. To avoid the exorbitant hurdles during the transition, the API must be non-intrusive in application code. On the other hand, the API needs to allow the internal implementation to timely \emph{intercept} signals to carry out communications and system optimizations. 
    \item \textbf{High Performance}: Data parallel training is subject to subtle dependencies between computations and communications. The design and implementation have to explore the solution space to efficiently convert more resources into higher training throughput.
\end{itemize}
PyTorch provides distributed data parallel as an \code{nn.Module} class, where applications provide their model at construction time as a sub-module. To guarantee mathematical equivalence, all replicas start from the same initial values for model parameters and synchronize gradients to keep parameters consistent across training iterations. To minimize the intrusiveness, the implementation exposes the same \code{forward}~\cite{forward} API as the user model, allowing applications to seamlessly replace subsequent occurrences of a user model with the distributed data parallel model object with no additional code changes. Several techniques are integrated into the design to deliver high-performance training, including bucketing gradients, overlapping communication with computation, and skipping synchronization.

Evaluations were conducted on an exclusive 32-GPU cluster and on 256 GPUs from a much larger shared entitlement. We developed benchmarks to evaluate the distributed package across different scales to present an in-depth view of the performance implications of different optimization techniques and configurations. Experiments also cover the comparison between NCCL and Gloo communication libraries. The results show that 1) communication is the dominant training latency contributor, and its impact increases with model sizes; 2) bucket sizes considerably affect communication efficiency, which could lead to more than 2X speedup if configured properly; 3) skipping synchronizations appropriately would significantly reduce amortized communication overhead without noticeably degrading convergence speed.


Techniques described in this paper were first released in PyTorch v1.1. During the past year, we have seen significant adoption both internally and externally. %
Within Facebook, a workload study from \code{05/11/20} to \code{06/05/20} shows that more than 60\% of production GPU hours during that period were spent on the PyTorch distributed data parallel package across a wide variety of applications, including speech, vision, mobile vision, translation, etc. 
%
%Within Facebook, more than 60\% of production GPU hours are spent on the PyTorch distributed data parallel package across a wide variety of applications, including speech, vision, mobile vision, translation, etc. 
There are three main contributions in this paper. First, this paper reveals the design and implementation of a widely adopted industrial state-of-the-art distributed training solution. Second, this paper highlights real-world caveats (\emph{e.g.}, due to pluralized graphs) that were overlooked by prior work. Third, we share performance tuning experiences collected from serving internal teams and open-source community users and summarized several directions for future improvements.


The remainder of the paper is organized as follows. Section~\ref{sec:background} briefly introduces PyTorch and data parallelism. Section~\ref{sec:design} elaborates the design for the PyTorch distributed data parallel module. Implementations and evaluations are presented in Section~\ref{sec:impl} and Section~\ref{sec:eval} respectively. Then, Section~\ref{sec:discussion} discusses lessons learned and opportunities for future improvements, and Section~\ref{sec:related} surveys related work. Finally, Section~\ref{sec:conclusion} concludes the paper.
